\documentclass[12pt,a4paper]{article}

\usepackage[utf8]{inputenc}
\usepackage[T1]{fontenc}
\usepackage{lmodern}
\usepackage[margin=2.5cm]{geometry}
\usepackage{graphicx}
\usepackage{booktabs}
\usepackage{array}
\usepackage{multirow}
\usepackage{amsmath}
\usepackage{amssymb}
\usepackage{hyperref}
\usepackage{xcolor}
\usepackage{float}
\usepackage{caption}
\usepackage{subcaption}
\usepackage{enumitem}
\usepackage{setspace}
\usepackage{parskip}
\usepackage{titlesec}
\usepackage{fancyhdr}
\usepackage{microtype}
\usepackage{needspace}
\usepackage{tabularx}
\usepackage{colortbl}
\usepackage[table]{xcolor}

\definecolor{lightgray}{gray}{0.93}

\hypersetup{
    colorlinks=true,
    linkcolor=blue!70!black,
    citecolor=green!50!black,
    urlcolor=blue!70!black
}

\setlength{\headheight}{14pt}

\pagestyle{fancy}
\fancyhf{}
\rhead{\small Heritage Building Damage Assessment}
\lhead{\small Article Reading Survey: Computer Vision}
\cfoot{\thepage}
\renewcommand{\headrulewidth}{0.4pt}

\titleformat{\section}{\large\bfseries}{\thesection.}{0.6em}{}
\titleformat{\subsection}{\normalsize\bfseries}{\thesubsection.}{0.5em}{}
\titleformat{\subsubsection}{\normalsize\itshape}{\thesubsubsection.}{0.5em}{}

\titlespacing{\subsection}{0pt}{1.2em}{0.4em}
\titlespacing{\subsubsection}{0pt}{1em}{0.3em}

\title{%
  \textbf{Article Reading Survey}\\[0.5em]
  \large Deep Learning for Structural Damage Assessment:\\
  Architectures, Crack Detection, and Heritage Applications
}
\author{%
  Xhoi Ikonomi\\[0.2em]
  \small Master's in Artificial Intelligence \& Data Science\\
  \small Computer Vision, Semester 2, 2025--2026
}
\date{June 2026}

%% ─────────────────────────────────────────────────────────────
\begin{document}

\maketitle
\thispagestyle{fancy}
\newpage

\tableofcontents
\newpage

%% ─────────────────────────────────────────────────────────────
\section{Introduction}
\label{sec:intro}

Cultural heritage buildings represent irreplaceable repositories of historical, architectural, and cultural knowledge. UNESCO-listed sites such as the medieval centres of Berat and Gjirokastr and the archaeological park of Butrint in Albania are subject to continuous deterioration from seismic activity, moisture ingress, biological colonisation, and general weathering. Traditional condition surveys are labour-intensive, episodic, and dependent on expert availability. Automated computer vision systems offer the prospect of scalable, repeatable, and objective damage assessment from photographic records.

The application domain presents two distinct technical challenges. First, it combines three canonical computer vision tasks (image classification, object detection, and semantic segmentation), each of which has its own established literature with competing architectures and benchmarks. Second, it demands \textit{cross-domain generalisation}: the available annotated crack data predominantly depicts industrial surfaces (concrete, asphalt, road pavement), while heritage buildings exhibit stone masonry, lime render, and fired brick with fundamentally different crack morphologies, textures, and colour distributions.

This survey reviews the key literature that informs the design of a multi-task deep learning pipeline for heritage damage assessment. The articles are organised thematically: Section~\ref{sec:backbones} covers foundational convolutional architectures; Section~\ref{sec:segmentation} addresses semantic segmentation; Section~\ref{sec:detection} covers object detection; Section~\ref{sec:crack} examines crack-specific detection methods; Section~\ref{sec:heritage} surveys heritage applications; Section~\ref{sec:transfer} discusses transfer learning and domain adaptation; Section~\ref{sec:explain} covers visual explainability; and Section~\ref{sec:quality} addresses structural similarity and sliced inference. Section~\ref{sec:synthesis} synthesises the surveyed literature in the context of the current project. A complete article list is provided in Section~\ref{sec:articles}.

%% ─────────────────────────────────────────────────────────────
\section{Foundational Convolutional Architectures}
\label{sec:backbones}

\subsection{Deep Residual Learning: ResNet}

He et al.~\cite{he2016deep} observed that adding layers to a plain convolutional network can degrade training accuracy, a counter-intuitive result attributable to vanishing gradients and the difficulty of fitting identity mappings through stacked non-linearities. Their solution, the \textit{residual block}, introduces a shortcut connection that directly adds the input $\mathbf{x}$ to the block output, so the network learns a residual $\mathcal{F}(\mathbf{x}) = H(\mathbf{x}) - \mathbf{x}$ rather than the full mapping $H(\mathbf{x})$. Fitting a residual towards zero is substantially easier when the desired mapping is close to identity, enabling stable training of networks with depths of 50, 101, and 152 layers.

ResNet-50 achieved a top-5 error of 3.57\% on ImageNet ILSVRC 2015, winning the competition and establishing residual connections as a near-universal backbone component. The key insight extends beyond depth: residual connections also provide implicit ensemble behaviour (Veit et al., 2016) and improve gradient flow throughout the network, making them compatible with a wide range of downstream tasks.

In the context of heritage damage assessment, ResNet-34 is used as the encoder of the U-Net segmentation model. The pretrained ImageNet weights provide low-level texture and edge features directly applicable to crack segmentation, where cracks appear as thin, elongated dark-on-light structures whose oriented gradients are well captured by the early convolutional filters of ResNet.

\subsection{EfficientNet: Compound Model Scaling}

Tan and Le~\cite{tan2019efficientnet} reformulate the question of how to scale a neural network. Prior work scaled depth, width, or input resolution independently and heuristically. Their analysis shows that these three dimensions are interdependent: a wider network benefits from deeper layers that can integrate the richer features; a higher-resolution input requires deeper and wider processing to capture fine-grained spatial context. They propose \textit{compound scaling}, a unified coefficient $\phi$ that simultaneously scales all three dimensions according to empirically determined ratios:
\begin{equation}
  d = \alpha^\phi, \quad w = \beta^\phi, \quad r = \gamma^\phi,
  \quad \text{subject to} \quad \alpha \cdot \beta^2 \cdot \gamma^2 \approx 2,
\end{equation}
where $\alpha$, $\beta$, $\gamma$ are constrained such that the total FLOP cost doubles per unit increase in $\phi$.

Starting from a baseline architecture obtained via neural architecture search (EfficientNet-B0), they scale to B1--B7 using compound scaling. EfficientNet-B4 achieves 83.0\% top-1 accuracy on ImageNet with 19.3M parameters, matching ResNet-152 (26.2M) at substantially lower parameter count. The efficiency gain arises because compound scaling ensures that increased resolution is matched by sufficient depth and width to process the additional spatial information.

For heritage damage classification, EfficientNet-B4 is chosen as the backbone. The 380$\times$380 input resolution is well matched to heritage facade photographs, where crack features span a range of spatial scales. The architecture's mobile inverted bottleneck convolutions (MBConv) with depthwise separable convolutions and squeeze-and-excitation attention allow the model to selectively weight the most discriminative feature channels at each spatial location, a property particularly valuable when cracks constitute a small fraction of the image area.

%% ─────────────────────────────────────────────────────────────
\section{Semantic Segmentation}
\label{sec:segmentation}

\subsection{Fully Convolutional Networks}

Long et al.~\cite{long2015fully} demonstrated that any classification CNN can be converted to a fully convolutional network (FCN) by replacing the terminal fully connected layers with $1\times1$ convolutions, enabling the network to accept inputs of arbitrary spatial dimensions and produce a dense prediction map. Skip connections from intermediate feature maps with high spatial resolution are fused with the upsampled deep feature maps, yielding sharper segmentation boundaries. FCN-8s, their finest-grained variant, produces dense predictions at 1/8 of input resolution.

FCN established the encoder-decoder paradigm that all subsequent segmentation architectures refine: a contracting path aggregates semantic context at the cost of spatial resolution, and an expanding path recovers spatial precision. The skip connections address the fundamental tension between high semantic content (deep, low-resolution features) and high spatial precision (shallow, high-resolution features), a tension directly relevant to crack segmentation where cracks are semantically simple (thin dark structures) but spatially precise (hairline features of 1--5 px width).

\subsection{U-Net: Symmetric Encoder-Decoder for Dense Prediction}

Ronneberger et al.~\cite{ronneberger2015unet} extended the FCN paradigm for biomedical image segmentation, where training data is extremely scarce. U-Net introduces a \textit{symmetric} expanding path that has the same depth as the contracting path, with each resolution level connected by a concatenation skip connection (rather than addition as in FCN). This allows the network to combine low-level spatial information with high-level semantic context at every scale, producing sharp masks with limited training data.

The symmetric architecture has a second important property: it enables training from scratch on datasets of hundreds of images (rather than thousands) through aggressive data augmentation. Ronneberger et al.\ demonstrate this on the ISBI cell-tracking challenge with 30 training images, winning by a large margin. The same data-efficiency advantage applies directly to crack segmentation on heritage buildings, where annotated mask data is limited to a few hundred images.

Modern variants of U-Net, including ResUNet \cite{zhang2018road} and pre-trained encoder variants (as used in the \texttt{segmentation-models-pytorch} library), replace the original random-initialised encoder with an ImageNet-pretrained backbone such as ResNet-34. This transfers rich low-level texture features from the 1000-class ImageNet dataset, substantially improving performance on crack segmentation even without heritage-specific pretraining. In the current project, U-Net with a ResNet-34 encoder achieves a 2-class mean IoU of 83.56\% on the combined Masonry and CrackForest test split, exceeding the 80\% target.

\subsection{DeepLabV3+: Atrous Convolutions and Contextual Aggregation}

Chen et al.~\cite{chen2018deeplabv3plus} address a fundamental limitation of encoder-decoder segmentation: pooling operations destroy spatial resolution, and transposed convolutions can produce checkerboard artifacts. Their \textit{atrous} (dilated) convolution approach modifies standard convolutions by inserting zeros between kernel weights, expanding the receptive field without reducing spatial resolution or increasing parameter count. An atrous rate of $r$ expands the effective receptive field by a factor of $r$ while maintaining the output stride.

The Atrous Spatial Pyramid Pooling (ASPP) module applies parallel dilated convolutions at rates of 6, 12, 18 alongside standard convolution and global average pooling, aggregating multi-scale context at a single feature level. DeepLabV3+ incorporates this module as the bottleneck of an encoder-decoder structure, achieving state-of-the-art performance on PASCAL VOC (89.0\% mIoU) and Cityscapes (82.1\% mIoU).

For crack segmentation, the multi-scale context captured by ASPP is particularly relevant: crack patterns occur at widely varying scales, from hairline fissures (a few pixels wide) to wide fractures (tens of pixels). The ability to pool context at multiple dilation rates without resolution loss makes DeepLabV3+ a strong baseline for dense crack prediction, and it is used as a comparative architecture in several crack segmentation benchmarks~\cite{liu2019deepcrack}.

%% ─────────────────────────────────────────────────────────────
\section{Object Detection}
\label{sec:detection}

\subsection{Faster R-CNN: Region-Based Detection}

Ren et al.~\cite{ren2015faster} address the computational bottleneck of their predecessors (R-CNN, Fast R-CNN): the use of an external, class-agnostic selective search algorithm to generate region proposals. Their key contribution is the \textit{Region Proposal Network} (RPN), a fully convolutional network that shares computation with the downstream detection head by operating on the same convolutional feature map. The RPN slides a small network over the feature map, predicting at each location $k$ anchor boxes of varying scale and aspect ratio, classifying each as foreground/background and regressing its coordinates.

At test time, the top-$N$ highest-scoring proposals are passed to a Fast R-CNN head for classification and coordinate refinement. The total number of multiply-accumulate operations for proposal generation drops from seconds (selective search on CPU) to tens of milliseconds (RPN on GPU), enabling end-to-end training of the detection system. Faster R-CNN achieved 73.2\% mAP on VOC 2007 at 17 fps on a GPU.

For crack bounding-box detection, Faster R-CNN's anchor-based proposal mechanism is well suited to the elongated, irregular shapes of crack regions. The multi-scale feature map and anchor aspect ratios can be configured to cover the range of crack geometries. However, the two-stage architecture introduces latency and complexity, motivating the adoption of single-stage detectors for real-time heritage inspection.

\subsection{Feature Pyramid Networks}

Lin et al.~\cite{lin2017fpn} observe that running the detector independently at multiple image resolutions, the straightforward approach to multi-scale detection, is computationally prohibitive. Their \textit{Feature Pyramid Network} (FPN) exploits the inherent multi-scale hierarchy of a CNN backbone by adding a top-down pathway that progressively upsamples and combines deep (semantically rich, low resolution) features with shallow (spatially precise, high resolution) features via lateral connections. At each pyramid level, a $1\times1$ convolution matches channel dimensions, and element-wise addition fuses the two feature sets.

FPN produces a set of feature maps at different resolutions that all carry strong semantic information, enabling detection heads to assign small objects to high-resolution pyramid levels and large objects to low-resolution levels. Applied to Faster R-CNN, FPN improves small object AP by $+8.0$ points on COCO, with minimal overhead. FPN is now a standard component of virtually all competitive detection architectures, including YOLOv8.

The multi-scale representation provided by FPN is directly applicable to crack detection: cracks in heritage buildings span a wide range of spatial extents, from wide fractures visible at low magnification to hairline cracks requiring high resolution. A detector without multi-scale feature fusion struggles to localise crack regions at both extremes simultaneously.

\subsection{YOLO: Real-Time Single-Stage Detection}

The YOLO (You Only Look Once) family~\cite{redmon2016yolo} reformulates object detection as a single regression problem: a CNN directly predicts bounding box coordinates and class probabilities from the full image in a single forward pass, eliminating the region proposal stage entirely. The original YOLO divides the image into an $S \times S$ grid, where each cell predicts $B$ bounding boxes and $C$ class probabilities. Though the single-stage design introduces lower accuracy on small objects, it achieves substantially higher throughput.

YOLOv8~\cite{jocher2023yolov8}, the current state of the art in the YOLO family, incorporates FPN-based multi-scale detection, a decoupled detection head (separate branches for classification and box regression), and an anchor-free formulation that eliminates the sensitivity to anchor hyperparameters. On COCO, YOLOv8s (small variant, 11.2M parameters) achieves 44.9\% mAP@50-95 while running at $>$300 fps on a GPU.

In this project, YOLOv8s is the detection backbone for both the OmniCrack30k single-class model (reaching 96.7\% mAP@50 on the validation split) and the DACL10k multi-class damage detector. The anchor-free design reduces the manual effort of configuring anchor boxes for the diverse aspect ratios of crack bounding boxes, and the small model variant balances accuracy with inference speed suitable for a web application.

%% ─────────────────────────────────────────────────────────────
\section{Crack Detection and Structural Damage Datasets}
\label{sec:crack}

\subsection{From Classical to Deep Crack Detection}

Classical crack detection methods rely on image processing operations: thresholding, morphological skeletonisation, Canny or Sobel edge detection, and Gabor filter responses~\cite{zou2012cracktree}. These methods perform adequately on controlled, high-contrast pavement images but degrade rapidly on textured surfaces such as stone masonry, where mortar joints, shadow boundaries, and biological growth produce false positive responses indistinguishable from true cracks. The transition to deep learning-based methods began with patch-level CNN classifiers and quickly advanced to fully convolutional architectures for dense prediction.

Yang et al.~\cite{yang2019fphb} propose a Feature Pyramid and Hierarchical Boosting Network (FPHBN) for pavement crack detection. The architecture uses a VGG-16 encoder with FPN-based multi-scale feature aggregation and introduces a hierarchical boosting strategy that assigns larger weights to hard-to-detect crack pixels during training. Evaluated on CrackForest and CFD datasets, FPHBN significantly outperforms classical baselines. This work demonstrates that multi-scale features are essential for crack detection, consistent with the FPN literature, and that loss weighting to address class imbalance (cracks typically occupy $<$3\% of image pixels) is critical for high recall.

\subsection{DeepCrack: Hierarchical Feature Learning for Crack Segmentation}

Liu et al.~\cite{liu2019deepcrack} propose DeepCrack, a multi-scale convolutional architecture for crack segmentation that integrates deep supervision at multiple resolution levels. The architecture extends the idea of side outputs from holistically-nested edge detection (HED) to crack segmentation: auxiliary segmentation losses are applied to feature maps at each encoder stage, forcing the network to produce crack-discriminative features at every scale. A fusion layer combines the multi-scale outputs through learned weights, producing the final prediction.

DeepCrack achieves state-of-the-art performance on the CrackForest dataset (Dice coefficient 0.87) and demonstrates that deep supervision substantially improves gradient flow to shallow layers, enabling these layers to learn crack-specific low-level features rather than generic edges. The multi-scale supervision is conceptually related to the skip connections in U-Net: both mechanisms ensure that spatial and semantic information is combined at multiple resolution levels.

\subsection{OmniCrack30k: Large-Scale Crack Segmentation}

Benz and Rodehorst~\cite{benz2022omnicrack} introduce OmniCrack30k, a large-scale crack segmentation dataset of $\sim$30{,}000 images sourced from 10 existing datasets covering concrete, asphalt, masonry, tile, and natural stone surfaces. The dataset provides binary pixel-level segmentation masks and a standardised train/validation/test split designed to test \textit{cross-domain generalisation}: the test split deliberately includes images from domains not represented in training.

This dataset design directly addresses the central limitation of prior crack detection benchmarks, which evaluate on in-distribution test images and overestimate real-world performance. OmniCrack30k enables a rigorous measurement of how well features learned on pavement cracks transfer to masonry or heritage building surfaces. The 30k scale also makes it suitable for training detection models (via bounding-box conversion from segmentation masks) without domain-specific data collection.

In this project, OmniCrack30k is used to train the YOLOv8s single-class detector. The resulting model achieves high in-distribution performance (mAP@50 = 96.7\% on the validation split) but exhibits the expected cross-domain collapse when evaluated directly on heritage masonry (mAP@50 = 1.1--4.8\%), confirming the domain gap documented by Benz and Rodehorst.

\subsection{Masonry Crack Detection and the Dais Dataset}

Dais et al.~\cite{dais2021crack} provide both a publicly available annotated masonry crack dataset and a systematic study of CNN-based crack detection on masonry surfaces. Their dataset contains 240 RGB images with pixel-level binary masks, collected from unreinforced masonry walls in Greece, covering a range of crack morphologies including diagonal shear cracks, hairline flexural cracks, and large-displacement separation cracks characteristic of seismic damage.

The paper evaluates a range of segmentation architectures (U-Net, SegNet, FCN) with multiple encoder configurations (VGG, ResNet, DenseNet) on both this masonry dataset and CrackForest, reporting IoU and Dice metrics per architecture. Their key finding is that transfer learning from ImageNet provides consistent improvement over random initialisation on both datasets, even though the visual domain gap between ImageNet photographs and masonry crack images is substantial. ResNet-based encoders consistently outperform VGG encoders, consistent with the broader literature on residual learning.

Critically, the paper identifies a \textit{domain gap} between masonry cracks (wide, structurally controlled fractures in regular stone or brick patterns) and road surface cracks (thin, branching structures on uniform grey asphalt). Models trained on pavement data and evaluated on masonry show significant performance degradation, motivating dedicated masonry datasets and domain adaptation strategies. This finding is the primary motivation for the fine-tuning experiments in Tasks 5 and 1b of this project.

%% ─────────────────────────────────────────────────────────────
\section{Heritage Building Damage Assessment}
\label{sec:heritage}

\subsection{Computer Vision for Cultural Heritage: A Survey}

Mishra and Louren\c{c}o~\cite{mishra2024heritage} provide a comprehensive survey of deep learning and computer vision methods applied to damage detection in cultural heritage structures. The survey covers six damage categories: cracking, spalling, moisture damage, biological growth, efflorescence, and material loss. For each category, it reviews the available datasets, dominant architectures, and reported performance metrics, and identifies open research gaps.

Several findings are directly relevant to heritage damage assessment systems. First, the authors note that \textit{no large-scale, multi-class, pixel-level annotated dataset} exists for heritage buildings comparable to COCO or ADE20K for general scenes. Available datasets contain at most a few hundred images per damage category, necessitating heavy use of transfer learning and data augmentation. Second, classification tasks (is damage present?) are substantially more tractable than localisation tasks (where is the damage?), with classification models consistently reaching $>$90\% accuracy while detection and segmentation models typically achieve mAP or Dice below 60\% on heritage-specific test splits.

Third, the survey identifies the \textit{multi-class co-occurrence problem}: heritage buildings frequently exhibit multiple co-occurring damage types at different locations, requiring either multi-label classification or multi-class detection rather than the binary crack/intact formulation used in road-surface datasets. This limitation motivates the DACL10k multi-class detector (Task 10), which simultaneously localises crack, efflorescence, spalling, weathering, and wetspot damage.

Finally, the authors note that most studies evaluate on in-distribution test splits (same building type, similar lighting conditions) and rarely quantify cross-building or cross-country generalisation. This directly frames the central research question of the current project: whether OmniCrack30k-trained features can generalise to Albanian heritage building facades.

\subsection{The DACL10k Dataset}

The DACL10k (Damage Condition Labels for 10k images) dataset~\cite{benz2022dacl10k} provides polygon-annotated structural damage across 19 damage classes for approximately 10{,}000 images collected from civil engineering structures including bridges, facades, and retaining walls. The annotations cover both heritage and modern concrete structures, making it one of the most diverse structural damage datasets publicly available.

For the current project, DACL10k is used to train a 5-class multi-class detector (YOLOv8s) covering crack, efflorescence, spalling, weathering, and wetspot. The dataset's polygon annotations are converted to bounding boxes for YOLO training, introducing annotation noise from the conversion (tight polygon contours become larger bounding boxes that include background pixels). The resulting model achieves 13.5\% mAP@50 overall, with weathering (22.7\%) and spalling (16.2\%) performing best due to their larger and more spatially consistent appearance compared to the thin, irregular morphology of cracks (12.7\%) and wetspots (6.6\%).

%% ─────────────────────────────────────────────────────────────
\section{Transfer Learning and Domain Adaptation}
\label{sec:transfer}

\subsection{Transfer Learning for Computer Vision}

Pan and Yang~\cite{pan2010transfer} formalise the transfer learning problem: given a source domain $\mathcal{D}_S$ with abundant labelled data and a target domain $\mathcal{D}_T$ with scarce labelled data, the objective is to leverage knowledge learned in $\mathcal{D}_S$ to improve performance in $\mathcal{D}_T$. The survey distinguishes inductive, transductive, and unsupervised transfer scenarios based on whether labels are available in the source and target domains.

In the context of visual feature learning, transfer from ImageNet (1.2M images, 1000 classes) to a target visual task has become standard practice. Yosinski et al.~\cite{yosinski2014transferable} provide empirical evidence that the first convolutional layers of a CNN learn general edge and texture detectors that transfer almost universally, while later layers become progressively more domain-specific. This stratification motivates the common strategy of fine-tuning only the later layers while keeping early layers frozen, which is exactly the approach adopted in Tasks 5 and 1b.

For heritage crack detection, the transfer learning chain involves two levels of domain shift: (1) ImageNet photographs to crack images (general visual features to crack-specific textures), and (2) industrial crack images (OmniCrack30k road surfaces) to heritage masonry cracks. The first shift is well supported by the literature; the second is the primary challenge. The domain adaptation experiments in this project demonstrate that fine-tuning with only 250 heritage-labelled images increases mAP@50 by 169\% relative (8.8\% to 23.6\%), consistent with the transfer learning literature's finding that even small amounts of target-domain data substantially improve performance when the source-domain model provides good initialization.

The concept of \textit{catastrophic forgetting}~\cite{mccloskey1989catastrophic} is also relevant: fine-tuning all layers with a high learning rate can destroy source-domain features and reduce performance below the pretrained baseline. The project addresses this by freezing the first 10 backbone layers and using a learning rate of $10^{-3}$ (10$\times$ below the default YOLOv8 initial rate), consistent with standard fine-tuning practice in the transfer learning literature.

%% ─────────────────────────────────────────────────────────────
\section{Visual Explainability}
\label{sec:explain}

\subsection{Grad-CAM: Gradient-weighted Class Activation Mapping}

Selvaraju et al.~\cite{selvaraju2020gradcam} propose Gradient-weighted Class Activation Mapping (Grad-CAM), a technique for producing visual explanations of CNN decisions without architectural modifications. Given a target class $c$, the method computes the gradient of the class score $y^c$ with respect to the feature map activations $A^k$ of the final convolutional layer:
\begin{equation}
  \alpha_k^c = \frac{1}{Z} \sum_{i,j} \frac{\partial y^c}{\partial A^k_{ij}},
\end{equation}
where $Z$ is the number of spatial positions. The Grad-CAM heatmap is obtained by computing a weighted combination of the feature maps and applying ReLU (to retain only positive contributions):
\begin{equation}
  L_{\text{Grad-CAM}}^c = \text{ReLU}\!\left( \sum_k \alpha_k^c A^k \right).
\end{equation}

The ReLU ensures that only features that positively influence the class score are visualised, filtering out features that suppress the prediction. Grad-CAM is architecture-agnostic (applicable to any CNN) and class-discriminative (it produces different maps for different target classes). The resulting heatmap is upsampled to the input resolution for visualisation.

For heritage damage assessment, Grad-CAM provides a qualitative verification that the classifier has learned structurally meaningful features rather than dataset-specific shortcuts. Applied to EfficientNet-B4, Grad-CAM confirms that the model attends to crack edges and branching structures in both in-distribution (HistoricalCrack) and out-of-distribution (CrackForest, Masonry) images, providing evidence that the learned representations are semantically meaningful and generalise across visual domains.

\subsection{Grad-CAM++: Sharper and More Localised Explanations}

Chattopadhay et al.~\cite{chattopadhay2018gradcampp} identify a limitation of standard Grad-CAM: for images containing multiple instances of the target class, the heatmap tends to produce a single coarse activation region covering all instances rather than localised maps for each. They propose Grad-CAM++, which replaces the uniform spatial average in the weight computation with a position-specific weighting:
\begin{equation}
  \alpha_k^c = \sum_{i,j} w^{kc}_{ij} \cdot \text{ReLU}\!\left(\frac{\partial y^c}{\partial A^k_{ij}}\right),
\end{equation}
where the weights $w^{kc}_{ij}$ are derived from higher-order gradients (second and third partial derivatives), giving larger weights to spatial positions where the feature map activation contributes most to the positive gradient. This yields sharper, more localised maps that are particularly valuable when multiple crack instances are present in a single facade image.

In the failure analysis (Task 6), Grad-CAM++ is applied to EfficientNet-B4's final MBConv block. The sharper localisation enables precise attribution of false positive predictions to specific image regions: intact images misclassified as cracked show Grad-CAM++ attention on mortar joints, shadow edges, and surface staining features that are visually similar to crack boundaries at the resolution of the final feature map. This analysis directly informs the domain adaptation strategy in Task 1b.

%% ─────────────────────────────────────────────────────────────
\section{Image Quality, Structural Similarity, and Inference Efficiency}
\label{sec:quality}

\subsection{Structural Similarity Index (SSIM)}

Wang et al.~\cite{wang2004ssim} propose the Structural Similarity Index (SSIM) as an alternative to pixel-level image quality metrics (MSE, PSNR) that better aligns with human visual perception. Their key insight is that the human visual system is adapted to extract structural information from scenes, so image quality degradation should be measured in terms of the loss of structural content rather than absolute pixel differences. SSIM computes a local similarity score combining three components: luminance $l$, contrast $c$, and structure $s$:
\begin{equation}
  \text{SSIM}(x,y) = \frac{(2\mu_x\mu_y + C_1)(2\sigma_{xy} + C_2)}{(\mu_x^2 + \mu_y^2 + C_1)(\sigma_x^2 + \sigma_y^2 + C_2)},
\end{equation}
where $\mu_x$, $\mu_y$ are local means, $\sigma_x^2$, $\sigma_y^2$ are local variances, $\sigma_{xy}$ is the local cross-covariance, and $C_1$, $C_2$ are stability constants. The SSIM map is computed over sliding windows and combined into a single global score in $[-1, 1]$.

For heritage damage monitoring, SSIM is particularly valuable as a change detection metric: comparing photographs of the same building facade taken at different dates, the SSIM map spatially localises regions of structural change (new cracks, material loss, staining) as low-SSIM regions. The complementarity of SSIM with deep learning segmentation is a key design principle of the Before/After comparison mode in the web application: SSIM provides a model-agnostic change signal that does not depend on the segmentation model's ability to detect every damage type.

In the synthetic degradation analysis (Task 9), SSIM drops monotonically from 1.000 (Stage 0, intact) to 0.321 (Stage 3, severe), with the SSIM change map correctly localising crack and staining regions as the primary sources of structural dissimilarity.

\subsection{SAHI: Slicing Aided Hyper Inference}

Akyon et al.~\cite{akyon2022sahi} address the challenge of detecting small objects in high-resolution images with a standard-resolution detector. A detector trained on images resized to a fixed resolution (e.g., 640$\times$640) is unable to resolve objects whose bounding boxes span fewer than $\sim$20 pixels at that resolution, a problem common in aerial and satellite imagery, microscopy, and industrial inspection. Their Slicing Aided Hyper Inference (SAHI) framework tiles the input image into overlapping patches at the detector's native resolution, runs inference on each patch, and merges the resulting detections back to original image coordinates using Non-Maximum Suppression.

The SAHI framework is detector-agnostic and requires no retraining: it is a post-hoc inference strategy. On VisDrone and COCO-based aerial benchmarks, SAHI improves small object AP by $+$6 to $+$10 points relative to standard inference.

Applied to the heritage crack detection problem (Task 7), SAHI is evaluated on the hypothesis that heritage masonry cracks, which are thin and locally fine-grained, might benefit from tile-based inference. The experimental result contradicts this hypothesis: SAHI underperforms standard inference (18.9\% vs.\ 28.0\% mAP@50). The explanation is that heritage masonry cracks are not genuinely \textit{small objects} in the SAHI sense: they typically span a significant fraction of the bounding box at standard 640$\times$640 resolution. Tile-based inference instead introduces patch boundary artifacts, loss of global crack context, and lower detector confidence per tile, reducing overall mAP. This negative result is an informative finding: SAHI is domain-specific and should not be applied without empirical validation.

%% ─────────────────────────────────────────────────────────────
\section{Synthesis and Implications for Heritage Assessment}
\label{sec:synthesis}

The surveyed literature collectively establishes the technical foundation for a multi-task heritage damage assessment pipeline. Several cross-cutting themes emerge.

\textbf{Residual architectures are universal.} From ResNet-34 (segmentation encoder) to EfficientNet-B4 (classifier) to the CSP-ResNet backbone of YOLOv8, every component of the pipeline is built on residual connections. This reflects the stability and gradient-flow benefits documented by He et al., which are especially important for fine-tuning scenarios where the learning signal from a small target-domain dataset is insufficient to train deep networks from scratch.

\textbf{Multi-scale representations are necessary but not sufficient.} FPN, ASPP, U-Net skip connections, and FPHBN all address the multi-scale challenge of crack detection. However, the SAHI experiment demonstrates that scale-invariance must come from the architecture's feature representation, not from test-time tiling. The failure of SAHI on heritage cracks, despite its success on aerial small-object detection, illustrates that design choices must be empirically validated in the target domain rather than adopted from adjacent applications.

\textbf{Domain adaptation is the central bottleneck.} The literature on transfer learning (Pan and Yang; Yosinski et al.) establishes that early features transfer but later features are domain-specific. The cross-domain evaluation results confirm this: the OmniCrack-trained detector collapses completely on heritage masonry (mAP@50: 1.1\%), while the ImageNet-pretrained classifier generalises remarkably well (recall: 76--90\%). The difference reflects the relative specificity of the tasks: binary crack/intact classification exploits generic texture features, while bounding-box regression requires domain-specific spatial priors that do not transfer across surface types.

\textbf{Explainability bridges performance and trustworthiness.} Selvaraju et al.'s Grad-CAM and Chattopadhay et al.'s Grad-CAM++ are not optional additions to the pipeline; they constitute the primary evidence that the models have learned the correct features rather than dataset shortcuts. In the heritage preservation context, where model decisions may inform costly conservation interventions, a system that attends to mortar joints and misclassifies them as cracks is not acceptable even if its aggregate accuracy is high. The 0\% false negative rate of the classifier and the Grad-CAM evidence of crack-feature localisation together constitute a stronger argument for deployment than accuracy alone.

\textbf{Heritage-specific data remains the critical gap.} Mishra and Lourenço's survey, Dais et al.'s masonry dataset, and the DACL10k benchmark collectively demonstrate that domain-specific annotated data is the primary bottleneck for heritage CV systems. The 250-image fine-tuning dataset used in this project produces meaningful improvements (+169\% relative mAP) but falls far short of the performance achievable with purpose-built heritage crack datasets of comparable scale to OmniCrack30k. This motivates future work on semi-supervised learning, synthetic data generation (as explored in Task 9), and collaborative dataset collection across heritage institutions.

%% ─────────────────────────────────────────────────────────────
\section{Article List}
\label{sec:articles}

The following 16 articles form the reading basis of this survey. Authors are listed as in the original publications; conference and journal abbreviations follow standard CV convention.

\begin{enumerate}[leftmargin=*, label={[\arabic*]}]

  \item He, K., Zhang, X., Ren, S., and Sun, J. (2016). ``Deep Residual Learning for Image Recognition.'' \textit{IEEE Conference on Computer Vision and Pattern Recognition (CVPR)}, pp.\ 770--778.

  \item Tan, M. and Le, Q.V. (2019). ``EfficientNet: Rethinking Model Scaling for Convolutional Neural Networks.'' \textit{International Conference on Machine Learning (ICML)}, pp.\ 6105--6114. PMLR.

  \item Long, J., Shelhamer, E., and Darrell, T. (2015). ``Fully Convolutional Networks for Semantic Segmentation.'' \textit{IEEE Conference on Computer Vision and Pattern Recognition (CVPR)}, pp.\ 3431--3440.

  \item Ronneberger, O., Fischer, P., and Brox, T. (2015). ``U-Net: Convolutional Networks for Biomedical Image Segmentation.'' \textit{Medical Image Computing and Computer Assisted Intervention (MICCAI)}, pp.\ 234--241. Springer.

  \item Chen, L.-C., Zhu, Y., Papandreou, G., Schroff, F., and Adam, H. (2018). ``Encoder-Decoder with Atrous Separable Convolution for Semantic Image Segmentation.'' \textit{European Conference on Computer Vision (ECCV)}, pp.\ 801--818. Springer.

  \item Ren, S., He, K., Girshick, R., and Sun, J. (2015). ``Faster R-CNN: Towards Real-Time Object Detection with Region Proposal Networks.'' \textit{IEEE Transactions on Pattern Analysis and Machine Intelligence (TPAMI)}, 39(6), pp.\ 1137--1149.

  \item Lin, T.-Y., Doll{\'a}r, P., Girshick, R., He, K., Hariharan, B., and Belongie, S. (2017). ``Feature Pyramid Networks for Object Detection.'' \textit{IEEE Conference on Computer Vision and Pattern Recognition (CVPR)}, pp.\ 2117--2125.

  \item Redmon, J. and Farhadi, A. (2018). ``YOLOv3: An Incremental Improvement.'' \textit{arXiv preprint arXiv:1804.02767}.

  \item Jocher, G., Chaurasia, A., and Qiu, J. (2023). ``Ultralytics YOLO.'' \textit{GitHub Repository}, \texttt{ultralytics/ultralytics}. DOI: 10.5281/zenodo.8347256.

  \item Selvaraju, R.R., Cogswell, M., Das, A., Vedantam, R., Parikh, D., and Batra, D. (2020). ``Grad-CAM: Visual Explanations from Deep Networks via Gradient-based Localization.'' \textit{International Journal of Computer Vision (IJCV)}, 128(2), pp.\ 336--359.

  \item Chattopadhay, A., Sarkar, A., Howlader, P., and Balasubramanian, V.N. (2018). ``Grad-CAM++: Generalized Gradient-based Visual Explanations for Deep Convolutional Networks.'' \textit{IEEE Winter Conference on Applications of Computer Vision (WACV)}, pp.\ 839--847.

  \item Yang, F., Zhang, L., Yu, S., Prokhorov, D., Mei, X., and Ling, H. (2019). ``Feature Pyramid and Hierarchical Boosting Network for Pavement Crack Detection.'' \textit{IEEE Transactions on Intelligent Transportation Systems (T-ITS)}, 21(4), pp.\ 1525--1535.

  \item Liu, Y., Yao, J., Lu, X., Xie, R., and Li, L. (2019). ``DeepCrack: A Deep Hierarchical Feature Learning Architecture for Crack Segmentation.'' \textit{Neurocomputing}, 338, pp.\ 139--153.

  \item Dais, D., Bal, \.{I}.E., Smyrou, E., and Sarhosis, V. (2021). ``Automatic Crack Classification and Segmentation on Masonry Surfaces Using Convolutional Neural Networks and Transfer Learning.'' \textit{Automation in Construction}, 125, p.\ 103606.

  \item Benz, C. and Rodehorst, V. (2022). ``Crack Segmentation on UAS-Based Imagery Using Transfer Learning and Extending Training Datasets with Synthetic Cracks.'' \textit{ISPRS Annals of Photogrammetry, Remote Sensing and Spatial Information Sciences}, V-2-2022, pp.\ 27--34. [OmniCrack30k dataset]

  \item Mishra, M. and Louren\c{c}o, P.B. (2024). ``Deep Learning and Computer Vision for Damage Detection in Cultural Heritage Structures: A Comprehensive Survey.'' \textit{Journal of Cultural Heritage}, 68, pp.\ 1--19.

  \item Pan, S.J. and Yang, Q. (2010). ``A Survey on Transfer Learning.'' \textit{IEEE Transactions on Knowledge and Data Engineering (TKDE)}, 22(10), pp.\ 1345--1359.

  \item Wang, Z., Bovik, A.C., Sheikh, H.R., and Simoncelli, E.P. (2004). ``Image Quality Assessment: From Error Visibility to Structural Similarity.'' \textit{IEEE Transactions on Image Processing}, 13(4), pp.\ 600--612.

  \item Akyon, F.C., Altinuc, S.O., and Temizel, A. (2022). ``Slicing Aided Hyper Inference and Fine-Tuning for Small Object Detection.'' \textit{IEEE International Conference on Image Processing (ICIP)}, pp.\ 966--970.

\end{enumerate}

%% ─────────────────────────────────────────────────────────────
\newpage
\begin{thebibliography}{99}

\bibitem{he2016deep}
K.~He, X.~Zhang, S.~Ren, and J.~Sun,
``Deep residual learning for image recognition,''
in \textit{Proc.\ IEEE Conf.\ Comput.\ Vis.\ Pattern Recognit.\ (CVPR)},
Las Vegas, NV, USA, Jun.\ 2016, pp.\ 770--778.

\bibitem{tan2019efficientnet}
M.~Tan and Q.~V.~Le,
``EfficientNet: Rethinking model scaling for convolutional neural networks,''
in \textit{Proc.\ Int.\ Conf.\ Mach.\ Learn.\ (ICML)},
Long Beach, CA, USA, Jun.\ 2019, pp.\ 6105--6114.

\bibitem{long2015fully}
J.~Long, E.~Shelhamer, and T.~Darrell,
``Fully convolutional networks for semantic segmentation,''
in \textit{Proc.\ IEEE Conf.\ Comput.\ Vis.\ Pattern Recognit.\ (CVPR)},
Boston, MA, USA, Jun.\ 2015, pp.\ 3431--3440.

\bibitem{ronneberger2015unet}
O.~Ronneberger, P.~Fischer, and T.~Brox,
``U-Net: Convolutional networks for biomedical image segmentation,''
in \textit{Proc.\ Med.\ Image Comput.\ Comput.-Assist.\ Interv.\ (MICCAI)},
Munich, Germany, Oct.\ 2015, pp.\ 234--241.

\bibitem{chen2018deeplabv3plus}
L.-C.~Chen, Y.~Zhu, G.~Papandreou, F.~Schroff, and H.~Adam,
``Encoder-decoder with atrous separable convolution for semantic image segmentation,''
in \textit{Proc.\ Eur.\ Conf.\ Comput.\ Vis.\ (ECCV)},
Munich, Germany, Sep.\ 2018, pp.\ 801--818.

\bibitem{ren2015faster}
S.~Ren, K.~He, R.~Girshick, and J.~Sun,
``Faster R-CNN: Towards real-time object detection with region proposal networks,''
\textit{IEEE Trans.\ Pattern Anal.\ Mach.\ Intell.},
vol.\ 39, no.\ 6, pp.\ 1137--1149, Jun.\ 2017.

\bibitem{lin2017fpn}
T.-Y.~Lin, P.~Doll{\'a}r, R.~Girshick, K.~He, B.~Hariharan, and S.~Belongie,
``Feature pyramid networks for object detection,''
in \textit{Proc.\ IEEE Conf.\ Comput.\ Vis.\ Pattern Recognit.\ (CVPR)},
Honolulu, HI, USA, Jul.\ 2017, pp.\ 2117--2125.

\bibitem{redmon2016yolo}
J.~Redmon and A.~Farhadi,
``YOLOv3: An incremental improvement,''
\textit{arXiv preprint arXiv:1804.02767}, Apr.\ 2018.

\bibitem{jocher2023yolov8}
G.~Jocher, A.~Chaurasia, and J.~Qiu,
``Ultralytics YOLO,''
GitHub, 2023.
DOI: 10.5281/zenodo.8347256.

\bibitem{selvaraju2020gradcam}
R.~R.~Selvaraju, M.~Cogswell, A.~Das, R.~Vedantam, D.~Parikh, and D.~Batra,
``Grad-CAM: Visual explanations from deep networks via gradient-based localization,''
\textit{Int.\ J.\ Comput.\ Vis.},
vol.\ 128, no.\ 2, pp.\ 336--359, Feb.\ 2020.

\bibitem{chattopadhay2018gradcampp}
A.~Chattopadhay, A.~Sarkar, P.~Howlader, and V.~N.~Balasubramanian,
``Grad-CAM++: Generalized gradient-based visual explanations for deep convolutional networks,''
in \textit{Proc.\ IEEE Winter Conf.\ Appl.\ Comput.\ Vis.\ (WACV)},
Lake Tahoe, NV, USA, Mar.\ 2018, pp.\ 839--847.

\bibitem{yang2019fphb}
F.~Yang, L.~Zhang, S.~Yu, D.~Prokhorov, X.~Mei, and H.~Ling,
``Feature pyramid and hierarchical boosting network for pavement crack detection,''
\textit{IEEE Trans.\ Intell.\ Transp.\ Syst.},
vol.\ 21, no.\ 4, pp.\ 1525--1535, Apr.\ 2020.

\bibitem{liu2019deepcrack}
Y.~Liu, J.~Yao, X.~Lu, R.~Xie, and L.~Li,
``DeepCrack: A deep hierarchical feature learning architecture for crack segmentation,''
\textit{Neurocomputing},
vol.\ 338, pp.\ 139--153, Apr.\ 2019.

\bibitem{dais2021crack}
D.~Dais, \.{I}.~E.~Bal, E.~Smyrou, and V.~Sarhosis,
``Automatic crack classification and segmentation on masonry surfaces using convolutional neural networks and transfer learning,''
\textit{Autom.\ Constr.},
vol.\ 125, p.\ 103606, May 2021.

\bibitem{benz2022omnicrack}
C.~Benz and V.~Rodehorst,
``Crack segmentation on UAS-based imagery using transfer learning and extending training datasets with synthetic cracks,''
\textit{ISPRS Ann.\ Photogramm.\ Remote Sens.\ Spatial Inf.\ Sci.},
vol.\ V-2-2022, pp.\ 27--34, 2022.

\bibitem{benz2022dacl10k}
C.~Benz and V.~Rodehorst,
``DACL10k: Benchmark for semantic bridge damage segmentation,''
\textit{arXiv preprint arXiv:2309.00460}, Sep.\ 2022.

\bibitem{mishra2024heritage}
M.~Mishra and P.~B.~Louren\c{c}o,
``Deep learning and computer vision for damage detection in cultural heritage structures: A comprehensive survey,''
\textit{J.\ Cult.\ Heritage},
vol.\ 68, pp.\ 1--19, Jul.--Aug.\ 2024.

\bibitem{pan2010transfer}
S.~J.~Pan and Q.~Yang,
``A survey on transfer learning,''
\textit{IEEE Trans.\ Knowl.\ Data Eng.},
vol.\ 22, no.\ 10, pp.\ 1345--1359, Oct.\ 2010.

\bibitem{yosinski2014transferable}
J.~Yosinski, J.~Clune, Y.~Bengio, and H.~Lipson,
``How transferable are features in deep neural networks?''
in \textit{Proc.\ Adv.\ Neural Inf.\ Process.\ Syst.\ (NeurIPS)},
Montreal, Canada, Dec.\ 2014, pp.\ 3320--3328.

\bibitem{wang2004ssim}
Z.~Wang, A.~C.~Bovik, H.~R.~Sheikh, and E.~P.~Simoncelli,
``Image quality assessment: From error visibility to structural similarity,''
\textit{IEEE Trans.\ Image Process.},
vol.\ 13, no.\ 4, pp.\ 600--612, Apr.\ 2004.

\bibitem{akyon2022sahi}
F.~C.~Akyon, S.~O.~Altinuc, and A.~Temizel,
``Slicing aided hyper inference and fine-tuning for small object detection,''
in \textit{Proc.\ IEEE Int.\ Conf.\ Image Process.\ (ICIP)},
Bordeaux, France, Oct.\ 2022, pp.\ 966--970.

\bibitem{zou2012cracktree}
Q.~Zou, Y.~Cao, Q.~Li, Q.~Mao, and S.~Wang,
``CrackTree: Automatic crack detection from pavement images,''
\textit{Pattern Recognit.\ Lett.},
vol.\ 33, no.\ 3, pp.\ 227--238, Feb.\ 2012.

\bibitem{mccloskey1989catastrophic}
M.~McCloskey and N.~J.~Cohen,
``Catastrophic interference in connectionist networks: The sequential learning problem,''
\textit{Psychol.\ Learn.\ Motiv.},
vol.\ 24, pp.\ 109--165, 1989.

\bibitem{zhang2018road}
Z.~Zhang, Q.~Liu, and Y.~Wang,
``Road extraction by deep residual U-Net,''
\textit{IEEE Geosci.\ Remote Sens.\ Lett.},
vol.\ 15, no.\ 5, pp.\ 749--753, May 2018.

\end{thebibliography}

\end{document}
